\chapter{Devenir masculin ou féminin}\label{ch:g11:becoming-male-female}

Sept semaines après la fécondation, un embryon humain possède une paire
de gonades qui ne sont ni des testicules ni des ovaires, et deux jeux
complets de canaux côte à côte --- l'un qui pourrait devenir le
tractus masculin, l'autre le tractus féminin. Rien de visible ne
distingue un embryon qui sera un garçon d'un embryon qui sera une
fille. Au fil des semaines suivantes, un jeu de canaux se développe et
l'autre se rétracte, les gonades prennent une forme ou l'autre, et au
quatrième mois l'anatomie est décidée. La décision est prise par un
seul \omterm{def:g10:universal-dna:gene}{gène}, et exécutée par deux hormones. Ce chapitre suit cette
succession, des \omterm{def:g10:universal-dna:information}{chromosomes} reçus à la fécondation jusqu'aux
changements de la \omterm{prop:g11:becoming-male-female:puberty}{puberté}.

\section{Les chromosomes sexuels}

\begin{definition}[Chromosomes sexuels]\label{def:g11:becoming-male-female:chromosomes}
Sur les 23 paires de \omterm{def:g10:universal-dna:information}{chromosomes} humains, 22 sont semblables dans les
deux sexes~; la 23\textsuperscript{e} paire est faite des
\emph{chromosomes sexuels}\index{chromosomes sexuels}~: deux
\omterm{def:g10:universal-dna:information}{chromosomes} X chez une femme (XX), un X et un Y chez un homme (XY). Le
Y est petit, porte peu de \omterm{def:g10:universal-dna:gene}{gènes} et ne se transmet que du père au fils~;
tout ovule porte un X, et un spermatozoïde porte soit un X soit un Y,
si bien que le spermatozoïde du père décide du sexe chromosomique de
l'enfant, à chances égales.
\end{definition}

\begin{proposition}[Le stade indifférencié]\label{prop:g11:becoming-male-female:undifferentiated}
Jusqu'à la septième semaine, les embryons des deux sexes chromosomiques
sont identiques par leur anatomie reproductrice~: une paire de gonades
indifférenciées, chacune flanquée de deux canaux --- le \emph{canal de
Wolff}, précurseur du tractus masculin (épididyme, canal déférent), et
le \emph{canal de Müller}, précurseur du tractus féminin (trompe,
utérus). Le développement sexuel est la transformation de ce plan
commun en l'une de deux formes~; c'est la différenciation de la gonade
qui l'enclenche.
\end{proposition}
\admitted

\begin{omfigure}
\begin{tikzpicture}[font=\small, >=Stealth, line cap=round]
  % undifferentiated stage
  \begin{scope}
    \draw[thick, fill=gray!25] (0,2.2) ellipse (0.35 and 0.55);
    \node[font=\scriptsize, anchor=south] at (0,2.85) {gonade};
    \draw[very thick, omDef] (-0.25,1.6) -- (-0.25,0);
    \draw[very thick, omThm] (0.25,1.6) -- (0.25,0);
    \node[font=\scriptsize, omDef, anchor=east] at (-0.35,0.8) {Wolff};
    \node[font=\scriptsize, omThm, anchor=west] at (0.35,0.45) {Müller};
    \node[anchor=north, align=center] at (0,-0.2) {semaine 7~:\\ les deux canaux};
  \end{scope}
  \draw[->, thick] (1.6,1.6) -- (3.0,3.4) node[midway, above left, font=\scriptsize] {XY};
  \draw[->, thick] (1.6,0.9) -- (3.0,-1.0) node[midway, below left, font=\scriptsize] {XX};
  % male
  \begin{scope}[xshift=4.2cm, yshift=2.0cm]
    \draw[thick, fill=omDef!45] (0,2.2) ellipse (0.35 and 0.55);
    \node[font=\scriptsize, anchor=south] at (0,2.85) {testicule};
    \draw[very thick, omDef] (-0.25,1.6) -- (-0.25,0);
    \draw[very thick, omThm, dotted] (0.25,1.6) -- (0.25,1.0);
    \node[font=\scriptsize, anchor=west, align=left] at (0.6,1.9) {testostérone~: le canal\\ de Wolff est maintenu et se développe};
    \node[font=\scriptsize, anchor=west, align=left] at (0.6,0.9) {AMH~: le canal de\\ Müller régresse};
    \node[anchor=north, align=center] at (0,-0.2) {tractus masculin};
  \end{scope}
  % female
  \begin{scope}[xshift=4.2cm, yshift=-2.8cm]
    \draw[thick, fill=omThm!35] (0,2.2) ellipse (0.35 and 0.55);
    \node[font=\scriptsize, anchor=south] at (0,2.85) {ovaire};
    \draw[very thick, omDef, dotted] (-0.25,1.6) -- (-0.25,1.0);
    \draw[very thick, omThm] (0.25,1.6) -- (0.25,0);
    \node[font=\scriptsize, anchor=west, align=left] at (0.6,1.9) {pas de testostérone~: le canal\\ de Wolff régresse};
    \node[font=\scriptsize, anchor=west, align=left] at (0.6,0.9) {pas d'AMH~: le canal de Müller\\ devient trompe et utérus};
    \node[anchor=north, align=center] at (0,-0.2) {tractus féminin};
  \end{scope}
\end{tikzpicture}

{\small Du plan commun à deux tractus. Le testicule, par ses deux
hormones, maintient le canal de Wolff et supprime celui de Müller~; en
l'absence de ces hormones, le canal de Müller se développe et celui de
Wolff régresse.}
\end{omfigure}

\section{Le gène qui décide}

\begin{proposition}[SRY et la gonade]\label{prop:g11:becoming-male-female:sry}
Le \omterm{def:g11:cell-cycle-mitosis:chromosome}{chromosome} Y porte un \omterm{def:g10:universal-dna:gene}{gène}, \emph{SRY}\index{gène SRY}, exprimé dans
la gonade indifférenciée à la septième semaine. Sa \omterm{def:g11:gene-expression:protein}{protéine} allume les
\omterm{def:g10:universal-dna:gene}{gènes} qui font de la gonade un \emph{testicule}. En son absence ---
chez un embryon XX --- la gonade devient un \emph{ovaire} quelques
semaines plus tard. C'est la présence de \omterm{prop:g11:becoming-male-female:sry}{SRY}, et non le nombre de
\omterm{def:g10:universal-dna:information}{chromosomes} X, qui décide de la gonade.
\end{proposition}
\begin{proof}[Faits expérimentaux]
De rares individus XY dont le \omterm{prop:g11:becoming-male-female:sry}{gène SRY} porte une \omterm{def:g11:mutations:mutation}{mutation} développent
des ovaires et un corps féminin~; de rares individus XX qui portent un
fragment du Y contenant \omterm{prop:g11:becoming-male-female:sry}{SRY}, transféré sur un X, développent des
testicules et un corps masculin. Chez la souris~: un embryon XX dans
lequel on a inséré le \omterm{def:g10:universal-dna:gene}{gène} \emph{Sry} de souris se développe en mâle.
Un seul \omterm{def:g10:universal-dna:gene}{gène} du Y suffit, et est nécessaire, pour faire un testicule.
\end{proof}

\begin{proposition}[Les hormones du testicule fœtal]\label{prop:g11:becoming-male-female:hormones}
Une fois formé, le testicule sécrète deux hormones. La
\emph{testostérone}\index{testostérone} maintient les canaux de Wolff
et les fait se développer en tractus interne masculin, et façonne les
organes génitaux externes sous la forme masculine. L'\emph{hormone
anti-müllérienne} (AMH) fait régresser les canaux de Müller. Un ovaire
n'en sécrète aucune à ce stade~: sans \omterm{prop:g11:becoming-male-female:hormones}{testostérone}, les canaux de Wolff
régressent, et sans AMH les canaux de Müller deviennent les trompes,
l'utérus et le haut du vagin. Le tractus féminin est donc la forme que
prennent les canaux en l'absence des hormones testiculaires.
\end{proposition}
\begin{proof}[Faits expérimentaux]
Jost (1947) a retiré les gonades de fœtus de lapin dans l'utérus
avant la différenciation~: tous les fœtus, XX ou XY, ont développé
un tractus féminin. Greffer un testicule à côté des gonades d'un
fœtus castré produisait un tractus masculin~; implanter un cristal
de \omterm{prop:g11:becoming-male-female:hormones}{testostérone} seule maintenait les canaux de Wolff mais ne
supprimait pas ceux de Müller, qui persistaient à côté d'eux. Le
testicule fabrique donc deux produits --- l'un qui bâtit les canaux
masculins (la \omterm{prop:g11:becoming-male-female:hormones}{testostérone}) et l'autre, identifié plus tard comme
l'AMH, qui supprime les canaux féminins --- et l'ovaire n'est pas
nécessaire à la formation du tractus féminin.
\end{proof}

\begin{omfigure}
\begin{tikzpicture}[font=\small, >=Stealth,
  b/.style={draw, rounded corners, align=center, text width=3.1cm, minimum height=1.1cm}]
  \node[b, fill=omDef!12] (a1) at (0,0) {fœtus, gonades retirées};
  \node[b, fill=omThm!12] (a2) at (4.2,0) {pas de testostérone,\\ pas d'AMH};
  \node[b, fill=white] (a3) at (8.4,0) {tractus féminin,\\ quels que soient les chromosomes};
  \draw[->] (a1) -- (a2); \draw[->] (a2) -- (a3);
  \node[b, fill=omDef!12] (b1) at (0,-1.7) {gonades retirées,\\ testicule greffé};
  \node[b, fill=omThm!12] (b2) at (4.2,-1.7) {testostérone\\ et AMH};
  \node[b, fill=white] (b3) at (8.4,-1.7) {tractus masculin};
  \draw[->] (b1) -- (b2); \draw[->] (b2) -- (b3);
  \node[b, fill=omDef!12] (c1) at (0,-3.4) {gonades retirées,\\ testostérone implantée};
  \node[b, fill=omThm!12] (c2) at (4.2,-3.4) {testostérone,\\ pas d'AMH};
  \node[b, fill=white] (c3) at (8.4,-3.4) {canaux masculins maintenus~;\\ canaux féminins persistants};
  \draw[->] (c1) -- (c2); \draw[->] (c2) -- (c3);
\end{tikzpicture}

{\small Les trois expériences de Jost sur des fœtus de lapin, et leur
lecture. Le tractus féminin n'a besoin d'aucune hormone~; le tractus
masculin a besoin de la \omterm{prop:g11:becoming-male-female:hormones}{testostérone} pour être bâti et d'une seconde
hormone testiculaire, l'AMH, pour supprimer les canaux féminins.}
\end{omfigure}

\begin{example}[La cascade dans l'ordre]\label{ex:g11:becoming-male-female:cascade}
Fécondation par un spermatozoïde \omterm{def:g11:genes-and-disease:genetic}{porteur} d'un Y~; semaine 7, \omterm{prop:g11:becoming-male-female:sry}{SRY}
s'exprime, la gonade devient un testicule~; semaines 8 à 12,
\omterm{prop:g11:becoming-male-female:hormones}{testostérone} et AMH~: canaux de Wolff maintenus, canaux de Müller
disparus, organes génitaux externes façonnés au masculin~; puis
quiescence jusqu'à la \omterm{prop:g11:becoming-male-female:puberty}{puberté}. Fécondation par un spermatozoïde \omterm{def:g11:genes-and-disease:genetic}{porteur}
d'un X~: pas de \omterm{prop:g11:becoming-male-female:sry}{SRY}~; semaines 10 à 12, la gonade devient un ovaire~;
pas d'hormones testiculaires~: les canaux de Müller se développent, les
canaux de Wolff régressent, les organes génitaux externes prennent la
forme féminine. Chaque étape dépend de la précédente, et un échec à
l'une quelconque change tout ce qui suit.
\end{example}

\section{La puberté}

\begin{proposition}[La puberté]\label{prop:g11:becoming-male-female:puberty}
Les organes reproducteurs, formés avant la naissance, restent inactifs
toute l'enfance. À la \emph{puberté}\index{puberté} --- vers onze à
treize ans, plus tôt chez les filles --- un centre du cerveau se met à
libérer, par impulsions, une hormone qui fait sécréter à l'hypophyse
deux hormones, LH et FSH, dans le sang. Celles-ci agissent sur les
gonades~: le testicule se met à fabriquer des spermatozoïdes et à
sécréter de la \omterm{prop:g11:becoming-male-female:hormones}{testostérone} en quantités adultes~; l'ovaire entame ses
cycles et sécrète des œstrogènes. Les hormones sexuelles produisent
alors les \emph{caractères sexuels secondaires}~: croissance des
organes reproducteurs, pilosité, voix, seins, répartition du muscle et
de la graisse, poussée de croissance puis soudure des \omterm{def:g10:muscles-and-joints:joint}{cartilages} de
croissance des os. La \omterm{prop:g11:becoming-male-female:puberty}{puberté} est achevée quand l'individu peut se
reproduire.
\end{proposition}
\admitted

\begin{omfigure}
\begin{tikzpicture}
\begin{axis}[
  width=0.86\linewidth, height=5.6cm,
  axis x line=bottom, axis y line=left,
  xmin=0, xmax=20, ymin=0, ymax=110,
  xlabel={âge (ans)}, ylabel={hormone sexuelle dans le sang (\% de l'adulte)},
  xlabel style={font=\small, at={(axis description cs:0.5,-0.12)}, anchor=north},
  ylabel style={font=\small},
  tick label style={font=\small}, xtick={0,4,8,12,16,20}, ytick={0,50,100},
  legend style={at={(0.03,0.95)}, anchor=north west, font=\small, draw=none},
]
\addplot[thick, omDef, smooth] coordinates {(0,3) (4,3) (8,3) (10,5) (11,12) (12,30) (13,55) (14,80) (15,95) (17,100) (20,100)};
\addplot[thick, omThm, smooth] coordinates {(0,3) (4,3) (8,4) (9,8) (10,20) (11,45) (12,70) (13,90) (15,100) (20,100)};
\legend{testostérone (garçons), œstrogènes (filles)}
\end{axis}
\end{tikzpicture}

{\small Les hormones sexuelles dans le sang au fil de l'enfance et de
l'adolescence, en pourcentage du niveau adulte (moyennes arrondies~; le
calendrier individuel varie de deux ou trois ans). Toutes deux sont
proches de zéro tant que le cerveau n'a pas donné le signal.}
\end{omfigure}

\begin{example}[Lire les courbes]\label{ex:g11:becoming-male-female:curves}
Les œstrogènes montent à partir de 9 ans environ chez les filles, la
\omterm{prop:g11:becoming-male-female:hormones}{testostérone} à partir de 11 ans chez les garçons~; chacune atteint la
moitié de son niveau adulte dans les deux ans qui suivent le départ, et
son plateau deux ou trois ans plus tard. La poussée de croissance suit
la montée d'environ un an, et les os cessent de grandir quand l'hormone
atteint le niveau adulte --- ce qui explique que les filles, qui
commencent plus tôt, cessent aussi de grandir plus tôt et soient en
moyenne moins grandes.
\end{example}

\section{Le sexe, en plusieurs sens}

\begin{proposition}[Distinguer les niveaux]\label{prop:g11:becoming-male-female:levels}
«~Masculin~» et «~féminin~» s'appliquent à plusieurs niveaux qui
s'accordent d'ordinaire mais restent distincts~:
\begin{itemize}
  \item le sexe \emph{chromosomique}~: XX ou XY, fixé à la
        fécondation~;
  \item le sexe \emph{gonadique}~: ovaire ou testicule, décidé par
        \omterm{prop:g11:becoming-male-female:sry}{SRY}~;
  \item le sexe \emph{anatomique}~: les tractus et les organes
        génitaux externes, façonnés par les hormones de la gonade
        fœtale, et les caractères secondaires, façonnés à la
        \omterm{prop:g11:becoming-male-female:puberty}{puberté}~;
  \item l'\emph{identité de genre}, le sentiment qu'a une personne
        d'être un homme, une femme ou ni l'un ni l'autre, et
        l'\emph{orientation sexuelle}, les personnes vers qui elle est
        attirée --- deux faits qui la concernent, qui ne se déduisent
        pas des niveaux précédents et qui ne se choisissent pas.
\end{itemize}
La biologie de ce chapitre décrit les trois premiers~; elle ne
détermine ni ne juge le dernier.
\end{proposition}
\admitted

\begin{example}[Quand les niveaux ne s'accordent pas]\label{ex:g11:becoming-male-female:disagree}
Environ une personne sur quelques milliers naît avec une combinaison
que la cascade ne produit pas d'ordinaire~: une personne XY dont les
\omterm{def:g10:cells-common-unit:cell}{cellules} n'ont pas le récepteur de la \omterm{prop:g11:becoming-male-female:hormones}{testostérone} développe un corps
féminin, avec des testicules jamais descendus et sans utérus~; une
personne XX porteuse d'une translocation de \omterm{prop:g11:becoming-male-female:sry}{SRY} qui forme un testicule
développe un corps masculin~; une personne XY dont le \omterm{prop:g11:becoming-male-female:sry}{SRY} ne fonctionne
pas se développe au féminin. Ces variations montrent, mieux que le cas
ordinaire, que chaque niveau est le résultat d'une étape précise --- un
\omterm{def:g10:universal-dna:gene}{gène}, une hormone, un récepteur --- et que les niveaux sont des faits
distincts.
\end{example}

\begin{method}[Raisonner sur un cas]\label{met:g11:becoming-male-female:case}
\begin{enumerate}
  \item Les \omterm{def:g10:universal-dna:information}{chromosomes}~: \omterm{prop:g11:becoming-male-female:sry}{SRY} est-il présent et fonctionnel~?
  \item La gonade~: testicule si \omterm{prop:g11:becoming-male-female:sry}{SRY} agit, ovaire sinon.
  \item Les hormones de la gonade fœtale~: \omterm{prop:g11:becoming-male-female:hormones}{testostérone} et AMH, ou
        ni l'une ni l'autre~; et les tissus peuvent-ils y répondre
        (récepteurs)~?
  \item Les canaux~: Wolff maintenu et Müller supprimé si les deux
        hormones agissent~; Müller développé si aucune n'agit~; les
        deux si la \omterm{prop:g11:becoming-male-female:hormones}{testostérone} agit sans l'AMH.
  \item La \omterm{prop:g11:becoming-male-female:puberty}{puberté}~: l'hormone sexuelle que sécrète la gonade décide
        des caractères secondaires --- de nouveau par l'intermédiaire
        des récepteurs.
\end{enumerate}
\end{method}

\begin{remark}[Un gène, et tout ce qui suit]\label{rem:g11:becoming-male-female:onegene}
Toute la différence entre les deux anatomies remonte à la présence ou
non d'un \omterm{def:g10:universal-dna:gene}{gène}, sur un \omterm{def:g11:cell-cycle-mitosis:chromosome}{chromosome}, à la septième semaine. Ce qui suit
est une cascade de biologie cellulaire ordinaire --- expression des
\omterm{def:g10:universal-dna:gene}{gènes}, hormones, récepteurs --- du genre que chaque chapitre de cette
année a décrit. Le plan dont part la cascade est commun aux deux sexes,
et la forme féminine est celle qu'elle bâtit quand le testicule
n'intervient pas.
\end{remark}

\section{Exercices}

\begin{exercise}[$\star$]\label{exo:g11:becoming-male-female:1}
Le gamète de quel parent décide du sexe chromosomique d'un enfant, et
pourquoi les chances sont-elles égales~?
\end{exercise}

\begin{exercise}[$\star$]\label{exo:g11:becoming-male-female:2}
Décrire l'anatomie reproductrice d'un embryon de sept semaines.
\end{exercise}

\begin{exercise}[$\star$]\label{exo:g11:becoming-male-female:3}
Que fait le \omterm{prop:g11:becoming-male-female:sry}{gène SRY}, et où se trouve-t-il~?
\end{exercise}

\begin{exercise}[$\star$]\label{exo:g11:becoming-male-female:4}
Nommer les deux hormones du testicule fœtal et donner l'effet de
chacune sur les canaux.
\end{exercise}

\begin{exercise}[$\star$]\label{exo:g11:becoming-male-female:5}
Qu'est-ce qui déclenche la \omterm{prop:g11:becoming-male-female:puberty}{puberté}, et quels sont les caractères
sexuels secondaires~?
\end{exercise}

\begin{exercise}[$\star\star$]\label{exo:g11:becoming-male-female:6}
Dans les expériences de Jost, qu'a montré le fœtus XY castré, et
qu'est-ce que cela a prouvé~?
\end{exercise}

\begin{exercise}[$\star\star$]\label{exo:g11:becoming-male-female:7}
Pourquoi l'implant de \omterm{prop:g11:becoming-male-female:hormones}{testostérone} seule a-t-il laissé les deux jeux de
canaux en place~? Qu'est-ce que cela a révélé~?
\end{exercise}

\begin{exercise}[$\star\star$]\label{exo:g11:becoming-male-female:8}
Un embryon de souris XX reçoit le \omterm{def:g10:universal-dna:gene}{gène} \emph{Sry} par \omterm{prop:g10:universal-dna:universal}{transgenèse}.
Prévoir ses gonades, ses canaux et son anatomie externe, et dire ce
qu'il ne pourra pas faire à l'âge adulte.
\end{exercise}

\begin{exercise}[$\star\star$]\label{exo:g11:becoming-male-female:9}
D'après la figure des hormones, à quel âge chaque hormone atteint-elle
la moitié de sa valeur adulte~? Quand chacune plafonne-t-elle~?
\end{exercise}

\begin{exercise}[$\star\star$]\label{exo:g11:becoming-male-female:10}
Expliquer, à l'aide des courbes, pourquoi les filles cessent de grandir
plus tôt que les garçons.
\end{exercise}

\begin{exercise}[$\star\star$]\label{exo:g11:becoming-male-female:11}
Un embryon XY porte une \omterm{def:g11:mutations:mutation}{mutation} qui détruit la \omterm{def:g11:gene-expression:protein}{protéine} \omterm{prop:g11:becoming-male-female:sry}{SRY}. Suivre la
cascade et décrire la personne à la naissance et à la \omterm{prop:g11:becoming-male-female:puberty}{puberté}.
\end{exercise}

\begin{exercise}[$\star\star\star$]\label{exo:g11:becoming-male-female:12}
Un fœtus XY a des testicules fonctionnels mais des \omterm{def:g10:cells-common-unit:cell}{cellules}
dépourvues du récepteur de la \omterm{prop:g11:becoming-male-female:hormones}{testostérone}~; l'AMH agit normalement.
Prévoir les gonades, chaque jeu de canaux, l'anatomie externe et les
caractères secondaires à la \omterm{prop:g11:becoming-male-female:puberty}{puberté} (les œstrogènes sont en partie
fabriqués à partir de la \omterm{prop:g11:becoming-male-female:hormones}{testostérone}).
\end{exercise}

\begin{exercise}[$\star\star\star$]\label{exo:g11:becoming-male-female:13}
Chez certains reptiles, c'est la température de l'œuf, et non un
\omterm{def:g11:cell-cycle-mitosis:chromosome}{chromosome}, qui décide du sexe. Expliquer en quel sens la cascade qui
suit la gonade pourrait être la même que chez les mammifères, et quelle
étape diffère.
\end{exercise}

\begin{exercise}[$\star\star\star$]\label{exo:g11:becoming-male-female:14}
Le \omterm{def:g11:cell-cycle-mitosis:chromosome}{chromosome} Y porte moins d'une centaine de \omterm{def:g10:universal-dna:gene}{gènes} et le X un millier
environ. Expliquer pourquoi une personne a besoin d'au moins un X mais
peut se passer d'un Y, et pourquoi les caractères liés au Y ne passent
que du père au fils.
\end{exercise}

\begin{exercise}[$\star\star\star$]\label{exo:g11:becoming-male-female:15}
Expliquer pourquoi l'affirmation «~la forme féminine est celle par
défaut du développement~» est une affirmation sur les hormones et non
sur l'importance, et ce qu'elle dit de l'évolution des deux sexes à
partir d'un plan commun.
\end{exercise}

\section{Problème~: les lapins de Jost}

\begin{problem}[{Problème du week-end --- les expériences qui ont
révélé les deux hormones du testicule, reconstituées étape par étape,
puis appliquées à trois cas humains et au calendrier de la puberté}]
\label{pb:g11:becoming-male-female:1}
Chez le lapin, les gonades se différencient vers le 19\textsuperscript{e}
jour de gestation et les canaux vers le 22\textsuperscript{e}. Jost a
opéré des fœtus dans l'utérus au 19\textsuperscript{e} jour et les a
examinés à la naissance (30\textsuperscript{e} jour).

\medskip
\noindent\textbf{Partie I --- Les expériences.}
\begin{enumerate}
  \item Fœtus des deux sexes chromosomiques, gonades retirées au
        19\textsuperscript{e} jour~: tous naissent avec des trompes et
        un utérus et sans canaux masculins. Que cela montre-t-il du
        rôle de l'ovaire, et de celui du testicule~?
  \item Pourquoi était-il indispensable d'opérer avant le
        22\textsuperscript{e} jour~?
  \item Fœtus castrés au 19\textsuperscript{e} jour, avec un
        testicule greffé dans l'abdomen~: canaux masculins, pas de
        canaux féminins. Que fournit la greffe, et comment atteint-elle
        les canaux~?
  \item Fœtus castrés au 19\textsuperscript{e} jour, avec un cristal
        de \omterm{prop:g11:becoming-male-female:hormones}{testostérone} implanté~: canaux masculins présents, canaux
        féminins présents aussi. Que fait la \omterm{prop:g11:becoming-male-female:hormones}{testostérone}, et que ne
        fait-elle pas~?
  \item En déduire que le testicule fabrique une seconde substance, et
        énoncer son effet. Comment son existence a-t-elle pu être
        confirmée ensuite~?
\end{enumerate}

\medskip
\noindent\textbf{Partie II --- Une seule gonade retirée.}
\begin{enumerate}[resume]
  \item On retire un seul testicule d'un fœtus XY au
        19\textsuperscript{e} jour. Prévoir les canaux de chaque côté
        du corps, sachant que les hormones agissent surtout du côté le
        plus proche.
  \item Un testicule est greffé à côté d'un ovaire d'un fœtus XX.
        Prévoir chaque côté.
  \item Un cristal de \omterm{prop:g11:becoming-male-female:hormones}{testostérone} est implanté à côté d'un ovaire d'un
        fœtus XX. Prévoir chaque côté.
  \item Pourquoi ces expériences unilatérales renforcent-elles la
        conclusion de la partie I~?
  \item Expliquer pourquoi un fœtus XX \omterm{def:g11:genes-and-disease:genetic}{porteur} d'une greffe de
        testicule serait néanmoins stérile à l'âge adulte.
\end{enumerate}

\medskip
\noindent\textbf{Partie III --- Trois cas humains.}
Appliquer la \cref{met:g11:becoming-male-female:case} à chacun.
\begin{enumerate}[resume]
  \item Une personne XX porte, sur l'un de ses \omterm{def:g10:universal-dna:information}{chromosomes} X, un
        fragment de Y contenant \omterm{prop:g11:becoming-male-female:sry}{SRY}.
  \item Une personne XY porte un \omterm{prop:g11:becoming-male-female:sry}{SRY} non fonctionnel.
  \item Une personne XY aux testicules normaux a des \omterm{def:g10:cells-common-unit:cell}{cellules}
        incapables de répondre à l'AMH.
  \item Dans lequel des trois cas la personne est-elle probablement
        fertile~? Justifier.
  \item Que montrent les trois cas au sujet du rapport entre sexe
        chromosomique et sexe anatomique~?
\end{enumerate}

\medskip
\noindent\textbf{Partie IV --- La \omterm{prop:g11:becoming-male-female:puberty}{puberté} en chiffres.}
D'après la figure des hormones, on admettra que la poussée de
croissance commence un an après que l'hormone a atteint 20~\% de son
niveau adulte et que les os cessent de grandir quand elle atteint
100~\%.
\begin{enumerate}[resume]
  \item Estimer l'âge auquel commence la poussée de croissance chez les
        filles et chez les garçons, et l'âge auquel la croissance
        s'arrête.
  \item Combien d'années de croissance chaque sexe a-t-il entre le
        début de la poussée et sa fin~?
  \item Un garçon grandit de \qty{8}{cm} par an pendant la poussée et
        une fille de \qty{7}{cm}. Estimer la taille gagnée pendant la
        poussée par chacun, et commenter l'écart moyen de taille
        adulte (environ \qty{13}{cm}).
  \item Une fille dont la montée des œstrogènes commence à 7 ans
        grandira vite tôt et s'arrêtera tôt. Expliquer pourquoi les
        médecins retardent parfois par des médicaments une \omterm{prop:g11:becoming-male-female:puberty}{puberté} très
        précoce.
  \item Énoncer le résultat~: les deux hormones que les lapins de Jost
        ont révélées, l'effet de chacune, et l'unique \omterm{def:g10:universal-dna:gene}{gène} qui allume
        l'organe qui les fabrique.
\end{enumerate}
\end{problem}
